\section{The tiered access path: from collection to admitted segment}
\label{sec:retrieval}
\label{sec:lanes}

The access path reads the collections, catalogue and indices maintained by
acquisition and organisation. It is the query-side counterpart of the
management levels: the catalogue tier chooses a collection, the shelf tier
chooses documents within it, the document tier opens a body, and the segment
tier can score its parts. Three access policies use different affordances of that
path, and a calling workflow may invoke the separate server-side Laya evidence
gate. The tool fetches referenced document text, segments it and applies a
local, query-conditioned verdict. Its model-card field (\texttt{noul}) is a
typed 0--1 score; this study reports no probability-calibration experiment and
does not interpret the score as a calibrated probability. Segments meeting the
configured threshold (0.5 by default) are returned as survivors with
provenance. A scored segment below the threshold is not admitted, while tool
or fetch failures are reported separately; a partially completed call may
still return successfully scored survivors. The answer workflow reserves a
scoped not-found report for a completed search with no supporting evidence,
not for an execution failure.

An access workflow can combine the following operations. Their order and
inclusion depend on the entry point and the question; the list is not a
server-enforced sequence for every query.
\begin{enumerate}
\item \emph{Rewrite.} The agent forms retrieval queries from the question.
\item \emph{Gather.} A pinned search calls two-stage or vector search; a
catalogue-first search selects collections and document descriptions. A
composed workflow can combine candidates from more than one source.
\item \emph{Merge and cap.} Where multiple sources are combined, the workflow
can deduplicate and limit the candidate set. Specific limits are properties of
the evaluated configuration, not universal platform guarantees.
\item \emph{Judge.} When invoked, the gate fetches candidate documents, scores
segments and returns survivors with provenance; fetch and tool failures are
reported separately from low scores.
\item \emph{Read and answer.} The agent examines the evidence pack and may read
a surviving document to check a specific claim before synthesising an answer.
\end{enumerate}

Figure~\ref{fig:tiered-algorithm} summarises this workflow and its terminal
conditions; Algorithm~\ref{alg:tiered} expresses the same decision sequence
in pseudocode. Both are workflow abstractions: individual entry points may
omit or combine candidate-gathering steps.

\begin{figure*}[t]
\centering
\includegraphics[width=\textwidth]{fig_tiered_algorithm}
\caption{Tiered access workflow. Candidate references can be gathered from
pinned collections, catalogue navigation or a combination. The optional Laya
gate fetches and scores segments. Incomplete review reports its scope and any
available evidence; a scoped not-found is reserved for a completed review
without support. The workflow describes intended agent behaviour, not a
server-enforced sequence.}
\label{fig:tiered-algorithm}
\end{figure*}

\begin{algorithm}[t]
\caption{Tiered access with optional evidence gate}
\label{alg:tiered}
\begin{algorithmic}[1]
\Require question $q$; catalogue $\mathcal{C}$; threshold $\tau$; gate flag $g$
\State $D \gets \textsc{GatherCandidates}(q,\mathcal{C})$ \Comment{pinned or catalogue-first}
\If{$g$}
  \State $S \gets \textsc{LayaJudge}(D,q,\tau)$ \Comment{fetch, score and return survivors}
\Else
  \State $S \gets \textsc{ReadSelected}(D)$
\EndIf
\If{execution incomplete}
  \State \Return incomplete status, searched scope, and available evidence
\ElsIf{completed review finds no support}
  \State \Return scoped not-found (searched scope, blind spots)
\Else
  \State \Return agent-synthesised answer with source references
\EndIf
\end{algorithmic}
\end{algorithm}
Exhausting a policy's call budget is intended to produce a scope and blind-spot
report. The endpoint enforces timeouts and turn limits; the per-policy budgets
are workflow instructions, not server quotas.

\subsection{Tier policies and the evidence gate}
\label{sec:trace}
\label{sec:gate}
\label{sec:gate-primitive}
\label{sec:twostage}

The current chat instructions provide two retrieval modes. In pinned mode,
the agent rewrites the question and calls the two-stage search (preferred) or
vector search within the selected collections; it can send candidate references
to Laya in sequential batches and optionally open a survivor for a specific
quote. In all-library mode, the instruction walks the catalogue and selected
shelves, selects document references from descriptions and sends these to Laya
before reading survivors. These are agent instructions, so actual tool use is
reported from traces rather than assumed from the instruction. Endpoint-level
timeouts and turn limits bound execution; retrieval sequences and batch sizes
are workflow instructions rather than separate server-enforced per-lane quotas.

The experiments also use fixed retrieval configurations. In particular, the
historical hybrid wrapper whose logged pools are analysed in
Section~\ref{sec:e6} is no longer a
registered tool. Another controlled candidate-pool experiment uses per-library
vector top-30 at cosine threshold 0.35 together with catalogue candidates,
followed by fusion and a twelve-candidate cap. These values identify that
experiment, not the current chat default. A composed hybrid directive may call
vector and catalogue tools in an agent workflow. Table~\ref{tab:lanes}
summarises the current modes and the evaluated composition.

The gate is the local Laya evidence classifier \cite{laya}, exposed as
\path{kb_laya_judge}. In reference mode, the server fetches candidate
documents, scores their segments against the query and returns surviving
segments with provenance and a bounded evidence pack. This is a threshold-based
admission operation; it does not itself fuse candidate sources or impose the
workflow's top-$k$ reading window. Those decisions belong to the calling search
configuration or agent workflow. To construct the evidence pack, the tool
orders surviving segments by source and position, deduplicates them and applies
the character limit. A passing verdict
is not proof of truth or entailment; the verifier remains subject to calibration
and transport limitations.

\begin{table*}[t]
\centering
\small
\setlength{\tabcolsep}{5pt}
\caption{Access policies over the tiered path. Modes differ in how candidate collections and documents are selected.
The hybrid row describes an experimental composition.}
\label{tab:lanes}
\begin{tabular}{@{}>{\raggedright\arraybackslash}p{0.11\textwidth}>{\raggedright\arraybackslash}p{0.19\textwidth}>{\raggedright\arraybackslash}p{0.24\textwidth}>{\raggedright\arraybackslash}p{0.16\textwidth}>{\raggedright\arraybackslash}p{0.20\textwidth}@{}}
\toprule
Policy & Tiers traversed & Retrieval affordance & Budget type & Intended terminal report\\
\midrule
Search & pinned collections, segment & Two-stage or vector search & Endpoint turn limit & Answer or report not-found\\
Librarian & catalogue, shelf, document, segment & Catalogue navigation and optional judging & Workflow instructions & Scope and blind spots\\
Hybrid & catalogue and segment & Composed vector and catalogue calls & Configuration-specific candidate cap & Candidates or judged survivors\\
\bottomrule
\end{tabular}
\end{table*}

The two-stage tool builds a candidate pool with lexical search and optional
graph expansion, then applies dense ranking. It does not call the
evidence gate; a ranked chunk is not admissible evidence until verified, and a
similarity threshold is not an abstention mechanism. The 13 October catalogue
incident showed that oversized tool returns can conceal documents before
reasoning begins. The tested repair adds a server-side character budget,
grouped views and pagination; its behaviour is reported in the Evaluation
section.

\subsection{Honest abstention}
A candidate list may contain high-similarity but unsupported documents. The
answer workflow instructs the agent to report not-found when its completed
evidence review finds no support, including the searched scope and known blind
spots, rather than falling back to parametric knowledge. This behaviour was
observed in the reported trials; it is not a server-side semantic-classifier
guarantee. Timeouts, unavailable tools and incomplete searches are execution
failures, not evidence of semantic absence (Section~\ref{sec:eval}).
